% Ներածություն (սևագիր v1)

\chapter*{Ներածություն}
\addcontentsline{toc}{chapter}{Ներածություն}

Ժամանակակից հաղորդակցական, համակարգչային, էներգետիկ և տրանսպորտային ցանցերը հնարավոր է մոդելավորել գրաֆներով, որտեղ գագաթները ցանցի հանգույցներն են, իսկ կողերը՝ նրանց միջև կապերը: Ցանցի հիմնական պահանջներից մեկը կապակցվածությունն է. հանգույցների միջև պետք է պահպանվի կապ նույնիսկ առանձին տարրերի խափանման դեպքում:

\textbf{Թեմայի արդիականությունը.} Ցանցերի աճող բարդության և կարևորության պատճառով անհրաժեշտ է ունենալ մեթոդներ, որոնք թույլ են տալիս քանակապես գնահատել ցանցի կայունությունը խափանումների նկատմամբ: Գրաֆների տեսությունը տալիս է կառուցվածքային չափանիշներ (կտրող կետեր, կամուրջներ, $k$-կապակցվածություն), իսկ հուսալիության տեսությունը՝ հավանականային գնահատականներ: Այդ երկու մոտեցումների համատեղ ուսումնասիրությունը թույլ է տալիս ոչ միայն գտնել ցանցի խոցելի տեղերը, այլև համեմատել տարբեր կառուցվածքները:

\textbf{Աշխատանքի նպատակը} գրաֆների կապակցվածության և ցանցի հուսալիության վերլուծության հիմնական մեթոդների ուսումնասիրությունն է, դրանց համեմատությունը և կիրառումը կոնկրետ ցանցերի վրա:

\textbf{Խնդիրները.}
\begin{enumerate}
  \item ներկայացնել գրաֆների կապակցվածության հիմնական հասկացությունները և թեորեմները.
  \item դիտարկել կապակցվածության որոշման և չափման ալգորիթմները.
  \item ներկայացնել ցանցի հուսալիության մոդելները և հաշվարկի մեթոդները (ճշգրիտ հաշվարկ, Monte Carlo սիմուլյացիա).
  \item համեմատել մեթոդները ըստ ճշգրտության և հաշվողական բարդության.
  \item իրականացնել ծրագրային մոդել և փորձարկել այն օրինակելի ցանցերի վրա:
\end{enumerate}

\textbf{Աշխատանքի կառուցվածքը.} Աշխատանքը բաղկացած է ներածությունից, երեք գլխից, եզրակացությունից և գրականության ցանկից: Առաջին գլխում ներկայացվում են տեսական հիմքերը, երկրորդում՝ մեթոդները, երրորդում՝ գործնական մասը:
